% =====================================================================
% 附录 A — 网盘文件说明
%
% 网盘随附文件清册。
% 列宽策略:文件名列定宽(\seqsplit 断词 + \texttt),说明列吃掉
% 全部剩余宽度;避免长文件名把说明列挤到重叠。
% =====================================================================

\chapter{网盘文件说明}\label{app:cloud-files}

为控制正文篇幅，部分大体积文件（安装包、系统架构图、演示视频等）不随正文装订，以网盘形式随作品提交。表~\ref{tab:cloud-files}列出网盘目录下的全部文件及其内容说明，供评审对照查阅。

\begin{table}[htbp]
\centering
\caption{网盘文件清单}\label{tab:cloud-files}
\begin{tabularx}{\linewidth}{@{}>{\raggedleft\arraybackslash}p{2.4em} l >{\raggedright\arraybackslash}X@{}}
\toprule
\textbf{序号} & \textbf{文件名} & \textbf{内容说明} \\
\midrule
1 & \seqsplit{BioMed-QAgent-1.0.0-win.zip} & Windows 安装包。 \\
2 & \seqsplit{BioMed-QAgent-1.0.0-x86\_64-Linux.AppImage} & Linux x86\_64 安装包。 \\
3 & \seqsplit{BioMed-QAgent-1.0.0.zip} & 纯构建产物，不包含运行环境。使用前需自行安装并配置 Python、uv 和 pnpm。 \\
4 & \seqsplit{biomed-qagent-main.pdf} & 系统架构图。 \\
5 & \seqsplit{qwen调用凭证.png} & Qwen 调用凭证。 \\
6 & \seqsplit{video.mp4} & 系统演示视频。 \\
\bottomrule
\end{tabularx}
\end{table}
